(1) 取时间因子 \(e^{-i\omega t}\)。质点处弦左右两段的横向张力合力为 \(T[y_x(L/2^+,t)-y_x(L/2^-,t)]\)，因此
\[
T[y_x(L/2^+,t)-y_x(L/2^-,t)]
=M\,y_{tt}(L/2,t),
\]
另有 \(y(L/2^-,t)=y(L/2^+,t)\)。该条件对横向力取实际代数方向，不能误把导数强行设为连续。

(2) 令 \(\ell=L/2\)、\(k=\omega/c\)。反对称模在中心有节点 \(Y(\ell)=0\)，左半区 \(Y=A\sin(kx)\)，故 \(\sin(k\ell)=0\)，即
\[
k_{\rm odd,n}=\frac{n\pi}{\ell}=\frac{2n\pi}{L},\quad n=1,2,\ldots.
\]
对称模左半区 \(Y=A\sin(kx)\)，右半区为 \(Y=A\sin[k(L-x)]\)。中心左右导数分别为 \(Ak\cos(k\ell)\) 和 \(-Ak\cos(k\ell)\)，代入跳跃条件：
\[
-2TkA\cos(k\ell)=-M\omega^2A\sin(k\ell).
\]
故对称模的非零频率满足
\[
2Tk\cos(k\ell)=M\omega^2\sin(k\ell),\quad
\cot(k\ell)=\frac{Mk}{2\mu}.
\]

(3) 反对称模质点恰不动，所受合力也因左右斜率相等而为零，因此频率与没有中心质点的相应节点模式相同。对一般模态，质点附加动能为 \(M|\dot y(\ell)|^2/2\)，故能量中的有效质量范数为
\[
\mathscr N=\mu\int_0^L|Y(x)|^2dx+M|Y(\ell)|^2>0.
\]
将波动方程乘以 \(Y^*\) 并在中心两侧分部积分，利用端点固定和质点跳跃条件，得到
\[
\omega^2\mathscr N=T\int_0^L|Y'(x)|^2dx>0.
\]
所有非零模均严格稳定。

(4) 当 \(M\to0^+\)，对称模的条件右端趋零，故最低对称根 \(k\ell\to\pi/2\)，频率 \(\omega_{\rm even,0}\to\pi c/L\)，恢复无点质量弦的最低模。当 \(M\to\infty\) 时最低根趋近零，利用 \(\cot\xi\sim1/\xi\)：
\[
\frac1{k\ell}\simeq\frac{Mk}{2\mu}
\quad\Longrightarrow\quad k^2\simeq\frac{2\mu}{M\ell}=\frac{4\mu}{ML}.
\]
故
\[
\omega_{\rm even,0}\sim 2\sqrt{\frac{T}{ML}}\to0.
\]
这是重质点缓慢移动、两半弦提供近似线性回复力的软模；若误以为 \(M\to\infty\) 时所有对称模都简单变成固定中点节点，会漏掉此低频分支。
